% Course report (PLAN Phase 6, item 2), derived from the paper: the same generated tables, numbers and
% results sentences (python -m ecgaug.paper), plus the end-to-end system the course asks for. Never type a
% number by hand; a red [bracket] is a result that does not exist yet.
\documentclass[11pt]{article}
\usepackage[margin=2.5cm]{geometry}
\usepackage{booktabs}
\usepackage{graphicx}
\usepackage{adjustbox}
\usepackage{amsmath}
\usepackage{xcolor}
\usepackage{hyperref}
\hypersetup{hidelinks}
\newcommand{\missing}[1]{\textcolor{red}{[#1]}}
% Generated by python -m ecgaug.paper; do not edit. Every value comes from a result file.
\newcommand{\CountTestF}{388}
\newcommand{\CountTestN}{44,259}
\newcommand{\CountTestS}{1,837}
\newcommand{\CountTestV}{3,221}
\newcommand{\CountTrainingF}{388}
\newcommand{\CountTrainingN}{35,496}
\newcommand{\CountTrainingS}{644}
\newcommand{\CountTrainingV}{3,030}
\newcommand{\CountValidationF}{27}
\newcommand{\CountValidationN}{10,370}
\newcommand{\CountValidationS}{300}
\newcommand{\CountValidationV}{758}
\newcommand{\PhaseOneInputRepresentationAMean}{0.164}
\newcommand{\PhaseOneInputRepresentationBMean}{0.001}
\newcommand{\PhaseOneInputRepresentationCMean}{0.296}
\newcommand{\PhaseOneInputRepresentationDMean}{0.035}
\newcommand{\PhaseOneWinner}{C}


\title{Rare-Arrhythmia Detection with Generative ECG Augmentation\\\large Course project report}
\author{\missing{team members}}
\date{}

\begin{document}
\maketitle

\section{Goal}
The system takes one ECG beat and returns a predicted arrhythmia class, a saliency map showing which parts of the
beat drove the prediction, and a short clinician-readable note grounded in the model output. Behind it is a research
question: under a strict inter-patient evaluation, does class-conditional generative augmentation (a conditional VAE
or a WGAN-GP) improve detection of the rare classes S and F compared with no augmentation, class-weighted loss,
random oversampling and SMOTE~\cite{dechazal2004,chawla2002smote,gulrajani2017wgangp,sohn2015cvae}?

\section{Data and protocol}
MIT-BIH Arrhythmia Database~\cite{moody2001,goldberger2000}, lead MLII, AAMI classes N, S, V and F, paced records
excluded. Training and validation come from the DS1 records and testing from the DS2 records~\cite{dechazal2004}, so no
patient is in two splits. The test set was read once, after every choice had been fixed on validation by rules written
down in advance.
\begin{table}[htbp]\centering\small% Generated by python -m ecgaug.paper from: results/phase1_data/record_counts.csv
\begin{tabular}{lrrrrr}
\toprule
Split & Records & N & S & V & F \\
\midrule
training & 17 & 35,496 & 644 & 3,030 & 388 \\
validation & 5 & 10,370 & 300 & 758 & 27 \\
test (DS2) & 22 & 44,259 & 1,837 & 3,221 & 388 \\
\bottomrule
\end{tabular}
\caption{Records and beats per split.}\end{table}

\section{System}
\begin{enumerate}
\item \textbf{Beat extraction}: a window around the reference R peak (representation \PhaseOneWinner, chosen on
validation), z-scored per beat, with RR features where the representation uses them.
\item \textbf{Classifier}: the cnn chosen at the phase 2 gate, trained with the augmentation condition chosen on
validation (\texttt{--condition best\_generative} by default).
\item \textbf{Saliency}: the method chosen per architecture at the phase 5 gate, summarised as the share of saliency in
the P, QRS, ST and T windows around the R peak.
\item \textbf{Note}: a template (or an LLM when configured) writes a note from a JSON object holding only the prediction,
the probabilities, a low-confidence flag and the top saliency segments; a fact checker rejects any claim not in that
object, and a template note is used whenever the check fails.
\end{enumerate}
Run it with \texttt{python -m ecgaug.demo --record <validation record> --beat <index>}; it prints the prediction,
the probabilities and the note, and saves a figure of the beat with its saliency.

\section{Methods in brief}
Four representations, four baseline classifiers, four quality-metric families, four cVAE and four WGAN-GP variants,
four augmentation policies, four saliency methods and four note generators were compared, each on validation and each
chosen by a rule fixed before the comparison. The final grid trains six conditions at three strengths for two
architectures and three seeds. Uncertainty on the test set comes from resampling the 22 test patients, with Holm
correction for the pre-registered comparisons~\cite{holm1979}. The paper (\texttt{paper/main.tex}) describes each
step in full.

\section{Results}
\begin{table}[htbp]\centering\small\adjustbox{max width=\linewidth}{% Generated by python -m ecgaug.paper from the result files
\begin{tabular}{lp{5.2cm}lp{5.2cm}}
\toprule
Hypothesis & Statement & Verdict & Evidence \\
\midrule
H1 & confirmatory: a generative condition improves test rare-class macro-F1 over none (Holm, 4 tests) & \missing{pending} &  \\
H1b & exploratory: gains larger for F than for S & \missing{pending} &  \\
H1c & secondary: generative beats SMOTE, random oversampling and class-weighted loss & \missing{pending} &  \\
H2 & exploratory: intra-patient evaluation inflates the baseline and the gains & \missing{pending} &  \\
H3 & exploratory: a quality metric predicts the downstream gain & \missing{pending} &  \\
H4 & exploratory: augmented models are at least as faithful as the baseline & \missing{pending} &  \\
\bottomrule
\end{tabular}
}\caption{Pre-registered hypotheses and their verdicts.}\end{table}
\missing{test results: written after /final-eval}

\begin{table}[htbp]\centering\small\adjustbox{max width=\linewidth}{\missing{test results: written by /final-eval}
}\caption{Test results per condition (95\% record-level CI).}\end{table}
\paragraph{Leakage.} \missing{leakage audit: optional, runs after the final evaluation}

\paragraph{Quality metrics.} \missing{RQ3: python -m ecgaug.eval.rq3 after the phase 3 runs}

\paragraph{Explanations.} \missing{RQ4: the phase 5 gate}

\paragraph{Notes.} \missing{note generators: the phase 5 gate}


\section{Limitations}
Twenty-two test patients, with the rare classes in few of them, give wide intervals; one database and one lead;
reference R peaks; synthetic beats get RR features copied from real beats; notes were rated by the team, not by
clinicians. The system is a course project and not a medical device.

\section{Team contributions}
\missing{who did what: from docs/ASSIGNMENTS.md once the draws are recorded}

\bibliographystyle{plain}
\bibliography{refs}
\end{document}
